% ---------------------------------------------------------------------------
% Part IX --- The physics beneath the circuits
% Written from Ezratty, Understanding Quantum Technologies 2026 (UQT 2026),
% chapters 3, 4.2, 4.4, 4.7, 5.1, 5.2, 6, 7.1, 7.2, 9.3, 10.5 and 12.1, and from
% the primary papers cited below.  Every figure of merit quoted with a date is a
% reported value, with its source; nothing here was measured by this repository.
% ---------------------------------------------------------------------------
\colorlet{cblue}{boxblue}
\colorlet{cverm}{boxred}
\definecolor{cpurple}{RGB}{104,46,130}
\definecolor{cgreen}{RGB}{32,110,60}
\definecolor{okfill}{RGB}{214,236,220}
\definecolor{lkfill}{RGB}{250,232,214}

\part*{Part IX --- The physics beneath the circuits}
\addcontentsline{toc}{part}{Part IX --- The physics beneath the circuits}

\begin{keybox}[title={What this part is for}]
Parts I--VIII count logical qubits and Toffolis, and \S\ref{rsa:sec:phys} converts them into
physical qubits and days with three numbers: a physical error rate $p=10^{-3}$, a $1\,\mu$s
surface-code cycle and a $10\,\mu$s reaction time. This part explains where those numbers come
from. It starts from the postulates and works down to the hardware: what a qubit is as a quantum
system, why a gate is a rotation and what physically produces it, what an error is and why a code
can discretise it, how transmons, ions, atoms, spins and photons each realise the two-level
system, what their measured fidelities and times were in September 2026, and how all of this feeds
the estimate. It ends with the record and the forecasts.

\medskip
The main source is Olivier Ezratty's \emph{Understanding Quantum Technologies 2026}
\cite{ph:uqt26}, a 1,800-page survey updated to September 2026, with the primary papers it rests
on cited where a number is quoted. The mathematics is worked out here rather than copied.
\end{keybox}

% ===========================================================================
\section{Quantum mechanics, as the circuits need it}\label{ph:sec:qm}

Everything in Parts I--VIII used four facts about quantum mechanics: states are vectors,
operations are unitary matrices, measurement is a projection with Born-rule probabilities, and
several systems compose by the tensor product. This section restates those facts with the care
needed to understand noise and error correction, and adds the two that the circuits never showed:
how a Hamiltonian turns into a gate, and what happens when the system is not isolated.

\subsection{States, observables, and the Born rule}

\paragraph{Pure states.} An isolated system is described by a unit vector $\ket\psi$ in a
complex Hilbert space $\mathcal H$. A qubit is any system for which we use a two-dimensional
subspace, $\ket\psi=\alpha\ket0+\beta\ket1$ with $|\alpha|^2+|\beta|^2=1$. A global phase is
unobservable, so a qubit has two real parameters, usually written
$\ket\psi=\cos\frac\theta2\ket0+e^{i\varphi}\sin\frac\theta2\ket1$: a point on the Bloch sphere,
with $\ket0$ at the north pole, $\ket1$ at the south, and the equator holding the equal
superpositions $\ket\pm=(\ket0\pm\ket1)/\sqrt2$ and $\ket{\pm i}=(\ket0\pm i\ket1)/\sqrt2$.

\paragraph{Observables.} A measurable quantity is a Hermitian operator $A=A^\dagger$. By the
spectral theorem $A=\sum_a a\,P_a$ with real eigenvalues $a$ and orthogonal projectors $P_a$
($P_a^2=P_a$, $P_aP_b=0$ for $a\ne b$, $\sum_aP_a=I$). The Pauli matrices
\[ X=\begin{pmatrix}0&1\\1&0\end{pmatrix},\quad
   Y=\begin{pmatrix}0&-i\\i&0\end{pmatrix},\quad
   Z=\begin{pmatrix}1&0\\0&-1\end{pmatrix} \]
are both Hermitian and unitary, square to $I$, anticommute pairwise and satisfy $XY=iZ$ and
cyclic permutations. Together with $I$ they span all $2\times2$ matrices: \emph{any}
single-qubit operator, in particular any error, is $aI+bX+cY+dZ$. That linear-algebra fact
is the whole reason quantum error correction is possible (\S\ref{ph:sec:open}).

\paragraph{Born rule and projection.} Measuring $A$ on $\ket\psi$ yields $a$ with probability
$p(a)=\bra\psi P_a\ket\psi=\|P_a\ket\psi\|^2$, after which the state is $P_a\ket\psi/\sqrt{p(a)}$.
Measuring again immediately returns the same $a$. For a qubit measured in the computational
basis, $p(0)=|\alpha|^2$ and $p(1)=|\beta|^2$; the expectation of $Z$ is $|\alpha|^2-|\beta|^2$.
Measuring in another basis means rotating first: an ``$X$ measurement'' is $H$ followed by a
$Z$ measurement, because $HZH=X$. This is what every stabiliser circuit in
\S\ref{ph:sec:codes} does, and what the ancilla measurements of \S\ref{im:sec:mbu}
did: an $X$-basis measurement is a Hadamard and a readout.

\paragraph{Mixed states.} When the state is not known exactly, or when we look at part of an
entangled system, the right object is the density operator $\rho=\sum_ip_i\ket{\psi_i}\bra{\psi_i}$:
positive, Hermitian, trace one, with $\langle A\rangle=\operatorname{tr}(A\rho)$ and
$p(a)=\operatorname{tr}(P_a\rho)$. A state is pure iff $\rho^2=\rho$, iff
$\operatorname{tr}\rho^2=1$. For one qubit every density operator is
\[ \rho=\tfrac12\big(I+r_xX+r_yY+r_zZ\big),\qquad |\vec r|\le1, \]
so mixed states fill the Bloch \emph{ball}; the surface is the pure states and the centre
$I/2$ is the maximally mixed state, which gives $50/50$ in every basis. Decoherence is the
motion of $\vec r$ into the ball (\S\ref{ph:sec:open}). The off-diagonal entries of $\rho$ in the
computational basis, $\rho_{01}=(r_x-ir_y)/2$, are the \emph{coherences}: they carry the relative
phase $\varphi$, and they are what dephasing destroys.

\subsection{Commutators, uncertainty, and why pulses must be slow}

Two observables can be known simultaneously iff they commute. If $[A,B]=AB-BA\ne0$ they cannot,
and the Robertson inequality quantifies it:
\[ \Delta A\,\Delta B\ \ge\ \tfrac12\big|\langle[A,B]\rangle\big| , \]
with $\Delta A^2=\langle A^2\rangle-\langle A\rangle^2$. For position and momentum $[x,p]=i\hbar$
gives Heisenberg's $\Delta x\,\Delta p\ge\hbar/2$. For a qubit, $[X,Z]=-2iY$, so
$\Delta X\,\Delta Z\ge|\langle Y\rangle|$: a state with a sharp $Z$ value has a random $X$ value,
which is exactly why measuring $X$ after $Z$ gives a coin flip.

The same mathematics governs control pulses, because a pulse of duration $\Delta t$ has a
frequency spread $\Delta f\gtrsim1/(4\pi\Delta t)$ (for a Gaussian envelope,
$\Delta t\cdot\Delta f\approx0.44$ in full widths). A $20$\,ns microwave pulse on a $5$\,GHz
transmon has a bandwidth of about $22$\,MHz. Two consequences run through
\S\ref{ph:sec:hw}: neighbouring qubits must be detuned from each other by more than that or the
pulse drives both (\emph{frequency crowding}), and a transmon's next level $\ket2$, only
$200$--$300$\,MHz away, is close enough that fast pulses leak into it (\emph{leakage}),
which is what DRAG pulse shaping is for.

\subsection{Unitary evolution: from Hamiltonian to gate}\label{ph:sec:unitary}

A closed system evolves by the Schr\"odinger equation
$i\hbar\,\partial_t\ket\psi=H\ket\psi$ with Hermitian Hamiltonian $H$. For time-independent
$H$ the solution is $\ket{\psi(t)}=U(t)\ket{\psi(0)}$ with
\[ U(t)=e^{-iHt/\hbar}, \]
unitary because $H$ is Hermitian. Every gate is such a $U$ for some $H$ applied for some time.
For a qubit, any Hamiltonian is $H=\tfrac{\hbar}{2}(\omega_0Z+\Omega\,\hat n\cdot\vec\sigma)$
up to a constant, and
\[ e^{-i\theta\,\hat n\cdot\vec\sigma/2}=\cos\tfrac\theta2\,I-i\sin\tfrac\theta2\,\hat n\cdot\vec\sigma \]
is the rotation of the Bloch vector by angle $\theta$ about the axis $\hat n$. So $X=-iR_x(\pi)$,
$H=-iR_y(\pi/2)R_x(\pi)$ up to phase, $S=R_z(\pi/2)$, $T=R_z(\pi/4)$: every single-qubit gate is a
rotation, and a rotation is a Hamiltonian switched on for a time. The small QFT rotations
$R_k=R_z(2\pi/2^k)$ of Part~I are physically the easiest gates of all --- a shorter pulse, or no
pulse at all --- which is why their fault-tolerant cost of $\approx60$--$100$ $T$ gates each
(Part~V) is so striking: the difficulty is not physics, it is that a code can protect only a
discrete set of rotations.

Two-qubit gates come from interaction Hamiltonians. Three examples cover the hardware:
\begin{itemize}\setlength\itemsep{2pt}
\item \textbf{$ZZ$ coupling.} $H=\tfrac{\hbar J}{4}Z\otimes Z$ for a time $t$ gives
  $U=\operatorname{diag}(e^{-i\phi},e^{i\phi},e^{i\phi},e^{-i\phi})$ with $\phi=Jt/4$. At
  $Jt=\pi$ this is $CZ$ up to single-qubit $Z$ rotations. Tunable-coupler transmons and
  Rydberg atoms both produce a state-dependent energy shift of this form.
\item \textbf{$XX$ coupling.} $H=\tfrac{\hbar\chi}{2}X\otimes X$ for $\chi t=\pi/2$ gives
  $e^{-i\frac\pi4X\otimes X}$, which is a maximally entangling gate equivalent to CNOT up to
  single-qubit rotations. This is the M\o lmer--S\o rensen gate of trapped ions.
\item \textbf{Exchange.} $H=J\,\vec S_1\cdot\vec S_2$ between two spins swaps them:
  $e^{-iHt/\hbar}$ at $Jt/\hbar=\pi$ is SWAP, at $\pi/2$ it is $\sqrt{\mathrm{SWAP}}$, and two
  $\sqrt{\mathrm{SWAP}}$s with a $Z$ rotation between them make a CNOT. This is the gate of
  silicon spin qubits.
\end{itemize}
In every case the entangling gate is slower and dirtier than single-qubit rotations because it
needs an interaction between two separately isolated objects, and the interaction must be
switched on and off cleanly. That is why two-qubit error rates dominate every error budget in
\S\ref{ph:sec:hw}.

\subsection{The driven two-level system}\label{ph:sec:rabi}

How is $R_x(\theta)$ actually produced? Take a qubit with splitting $\hbar\omega_0$ and drive it
with an oscillating field at frequency $\omega$, amplitude $\Omega$ and phase $\phi$:
\[ H=\tfrac{\hbar\omega_0}{2}Z+\hbar\Omega\cos(\omega t+\phi)\,X . \]
Go to the frame rotating with the drive, $\ket\psi\to e^{i\omega tZ/2}\ket\psi$, and drop the terms
oscillating at $2\omega$ (the rotating-wave approximation, valid for $\Omega\ll\omega$):
\[ H'=\tfrac{\hbar\Delta}{2}Z+\tfrac{\hbar\Omega}{2}\big(\cos\phi\,X+\sin\phi\,Y\big),
   \qquad \Delta=\omega_0-\omega . \]
On resonance ($\Delta=0$) the qubit rotates about an equatorial axis at angle $\phi$, at the Rabi
frequency $\Omega$: the population of $\ket1$ oscillates as $\sin^2(\Omega t/2)$. A pulse of area
$\Omega t=\pi$ is an $X$ gate, $\pi/2$ is a half-turn to the equator, and the \emph{phase} of
the drive chooses the axis, so $Y$ rotations cost nothing extra. A $Z$ rotation needs no pulse
at all: shifting the phase of all later pulses by $\theta$ is the same as $R_z(\theta)$, the
\emph{virtual $Z$} that every control stack uses. Off resonance the effective axis tilts out of
the equator and the oscillation amplitude drops to $\Omega^2/(\Omega^2+\Delta^2)$, which is how a
drive addresses one qubit and not its neighbour at a different frequency.

Numbers make this concrete. A transmon has $\omega_0/2\pi\approx5$\,GHz and
$\Omega/2\pi\approx10$--$50$\,MHz, so a $\pi$ pulse takes $10$--$50$\,ns and contains $50$--$250$
carrier oscillations. A hyperfine ion qubit has $\omega_0/2\pi\approx3$--$13$\,GHz driven by
microwaves or Raman lasers with $\Omega/2\pi\approx10$--$100$\,kHz, so a $\pi$ pulse takes
$5$--$50\,\mu$s. The gate time is set by the drive strength the hardware tolerates, not by the
qubit frequency: faster drives need more bandwidth (previous subsection), more power, and leak.

\subsection{Composite systems, entanglement, and two theorems}\label{ph:sec:ent}

Two systems have state space $\mathcal H_A\otimes\mathcal H_B$. A product state
$\ket a\otimes\ket b$ describes two independent objects; a state that cannot be written as a
product is \emph{entangled}. The Bell state $\ket{\Phi^+}=(\ket{00}+\ket{11})/\sqrt2$ is the
standard example: its reduced state $\rho_A=\operatorname{tr}_B\ket{\Phi^+}\bra{\Phi^+}=I/2$ is
maximally mixed, so each half alone is a fair coin while the two halves are perfectly correlated
in \emph{every} basis ($\ket{\Phi^+}=(\ket{++}+\ket{--})/\sqrt2$ too). The bilinearity of
$\otimes$ used in Part~I to explain phase kickback is the same structure: $(\alpha\ket0+\beta\ket1)\otimes\ket\phi$
is a sum of two product terms, and a controlled gate acts on each term separately.

\paragraph{CHSH.} That these correlations are not classical is a theorem with a two-line
calculation. Let Alice measure $A_0$ or $A_1$ and Bob $B_0$ or $B_1$, each with outcomes $\pm1$.
If outcomes were fixed in advance (realism) and independent of the other side's choice
(locality), then for every run $A_0B_0+A_0B_1+A_1B_0-A_1B_1=A_0(B_0+B_1)+A_1(B_0-B_1)=\pm2$, so
\[ S=\langle A_0B_0\rangle+\langle A_0B_1\rangle+\langle A_1B_0\rangle-\langle A_1B_1\rangle\le2 . \]
For $\ket{\Phi^+}$ and spin measurements along directions $\hat a,\hat b$ in the $x$--$z$ plane,
$\langle(\hat a\cdot\vec\sigma)\otimes(\hat b\cdot\vec\sigma)\rangle=\cos(a-b)$. Choosing
$a_0=0$, $a_1=90^\circ$, $b_0=45^\circ$, $b_1=135^\circ$ gives
$S=3\cos45^\circ-\cos135^\circ=2\sqrt2\approx2.83$, Tsirelson's bound. Experiments since 1982
find $S>2$; the loophole-free ones of 2015 (NV centres at Delft, photons at Vienna and NIST)
closed the detection and locality loopholes at once \cite{ph:uqt26}. The $2\sqrt2$ is the
maximum any quantum state can give, so a measured $S$ is also a diagnostic: it is the figure of
merit for two-qubit entanglement on a new processor.

\paragraph{No-cloning.} Suppose a unitary $U$ copies unknown states:
$U\ket\psi\ket0=\ket\psi\ket\psi$ and $U\ket\phi\ket0=\ket\phi\ket\phi$. Unitaries preserve inner
products, so $\langle\phi|\psi\rangle=\langle\phi|\psi\rangle^2$, which forces
$\langle\phi|\psi\rangle\in\{0,1\}$: only orthogonal states can be copied. Equivalently, by
linearity, $U(\alpha\ket0+\beta\ket1)\ket0=\alpha\ket{00}+\beta\ket{11}$, an entangled state, not
$(\alpha\ket0+\beta\ket1)^{\otimes2}$, which has cross terms. Three consequences shape the
rest of this part: a code cannot protect a qubit by keeping copies and voting, so it must
spread the state over many qubits and measure only \emph{parities} (\S\ref{ph:sec:codes}); a
magic state cannot be duplicated, so every $T$ gate consumes a freshly made one
(\S\ref{ph:sec:ft}); and the only way to move a state is to \emph{teleport} it, which is
precisely how a $T$ gate is applied to a logical qubit.

\subsection{Open systems: what an error is}\label{ph:sec:open}

A qubit is never isolated. Coupling to an environment that we do not track turns the pure-state
dynamics into dynamics of $\rho$. When the environment has no memory on the time scale of the
qubit (Markovian noise) the evolution is the Lindblad equation
\[ \frac{d\rho}{dt}=-\frac i\hbar[H,\rho]+\sum_k\Big(L_k\rho L_k^\dagger-\tfrac12\{L_k^\dagger L_k,\rho\}\Big), \]
whose jump operators $L_k$ encode the noise channels. Two of them describe nearly everything
measured on a qubit:
\begin{itemize}\setlength\itemsep{2pt}
\item \textbf{Energy relaxation}, $L=\sqrt{\gamma_1}\,\sigma_-$ with $\sigma_-=\ket0\bra1$: the
  excited population decays as $e^{-t/T_1}$, $T_1=1/\gamma_1$. It is a loss of energy to the
  environment (a photon into a lossy dielectric, a spontaneously emitted photon, a quasiparticle).
  It is \emph{not} a bit flip: it is non-unital, it pushes every state towards $\ket0$, and it
  also damps the coherence at rate $\gamma_1/2$.
\item \textbf{Pure dephasing}, $L=\sqrt{\gamma_\varphi/2}\,Z$: the populations are untouched and
  the coherences decay as $e^{-t/T_\varphi}$. Physically the qubit frequency fluctuates, so the
  phase $\varphi$ accumulated in a superposition becomes random.
\end{itemize}
Together, the coherence decays as $e^{-t/T_2}$ with
\[ \frac1{T_2}=\frac1{2T_1}+\frac1{T_\varphi},\qquad\text{so}\qquad T_2\le2T_1 . \]
$T_1$ is measured by preparing $\ket1$ and watching it decay; $T_2^*$ by a Ramsey sequence
($\pi/2$ pulse, wait $t$, $\pi/2$ pulse, read out: the fringes decay), and $T_2$ by a Hahn echo
that inserts a $\pi$ pulse at $t/2$ to refocus slow frequency drifts, so $T_2^*\le T_2$. All three
are fitted envelope constants, not the lifetime of ``the superposition''.

\paragraph{Channels and Kraus operators.} Integrated over a gate time, a noise process is a
completely positive trace-preserving map $\rho\mapsto\sum_kE_k\rho E_k^\dagger$ with
$\sum_kE_k^\dagger E_k=I$. Amplitude damping with decay probability $\gamma$ has
\[ E_0=\begin{pmatrix}1&0\\0&\sqrt{1-\gamma}\end{pmatrix},\qquad
   E_1=\begin{pmatrix}0&\sqrt\gamma\\0&0\end{pmatrix}; \]
pure dephasing is $\rho\mapsto(1-p)\rho+pZ\rho Z$; the depolarising channel is
$\rho\mapsto(1-p)\rho+\tfrac p3(X\rho X+Y\rho Y+Z\rho Z)$, which contracts the Bloch vector
uniformly towards the origin. A \emph{coherent} error is different in kind: it is a small unitary
over-rotation $e^{-i\epsilon\hat n\cdot\vec\sigma/2}$ from a miscalibrated pulse, and it adds in
amplitude over repeated gates (error $\propto\epsilon^2m^2$ after $m$ gates) where stochastic
errors add in probability ($\propto pm$). Error correction and randomised benchmarking both tend
to convert coherent errors into stochastic ones, which is one reason the two are budgeted
differently.

\paragraph{Why a code can correct a continuum of errors.} An arbitrary single-qubit error is
$E=aI+bX+cY+dZ$ with complex coefficients, a continuum. A stabiliser code measures operators
that commute or anticommute with each Pauli, so after the syndrome measurement the state is
\emph{projected} onto one of the four branches: no error, $X$, $Y$ or $Z$, with probabilities
$|a|^2,|b|^2,|c|^2,|d|^2$. The measurement discretises the error, and the correction undoes a
Pauli. A small over-rotation $e^{-i\epsilon X/2}=\cos\frac\epsilon2I-i\sin\frac\epsilon2X$
therefore becomes ``no error'' with probability $\cos^2\frac\epsilon2$ and a clean $X$ flip with
probability $\sin^2\frac\epsilon2$. This is the content of the Knill--Laflamme conditions: it is
enough to correct $X$, $Y=iXZ$ and $Z$ on every qubit. It is also why leakage out of the
$\{\ket0,\ket1\}$ subspace is a different and nastier beast: it is not spanned by the Paulis.

\subsection{Measurement, fidelity, and the error budget}\label{ph:sec:fid}

\paragraph{Projective and dispersive readout.} Textbook measurement projects onto $\ket0$ or
$\ket1$ and leaves the qubit there. Real readout couples the qubit to a meter --- a resonator, a
fluorescence cycle, a charge sensor --- and reads the meter. A readout is \emph{quantum
non-demolition} (QND) if it leaves the eigenstate it reports in place, so it can be repeated;
error correction needs this for the ancillas, which are read every cycle. The readout error
(``assignment error'') is the probability of reporting the wrong basis state; combined with
preparation errors it is quoted as SPAM.

\paragraph{Fidelity and randomised benchmarking.} The fidelity of a gate is the overlap between
what it does and what it should do, averaged over inputs; the \emph{error} is $1-F$. Fidelities
are measured by randomised benchmarking: run random sequences of $m$ Clifford gates followed by
their inverse, measure the survival probability of $\ket0$, and fit $A\,p^m+B$. Random Cliffords
twirl any noise into a depolarising channel, so the decay constant gives the average error per
Clifford
\[ r=(1-p)\,\frac{d-1}{d},\qquad r=\tfrac12(1-p)\ \text{for one qubit},\quad r=\tfrac34(1-p)\ \text{for two}. \]
Interleaving a specific gate between the random ones isolates that gate's error. Cross-entropy
benchmarking (XEB) does the same for whole random circuits and was the metric of the 2019
supremacy experiment. All of these report \emph{average} errors; the tails, crosstalk and
leakage that matter for a code are only partly visible in them.

\paragraph{The budget.} A circuit with $N$ qubits and depth $D$ has about $ND$ locations at which
an error can occur. With independent errors of probability $\epsilon$ per location, the expected
number of faults is $ND\epsilon$, and the circuit has a fair chance of running cleanly only if
\[ N\,D\,\epsilon\lesssim1 . \]
This is the heuristic behind every number in \S\ref{rsa:sec:phys}. The Shor circuits of Parts
IV--VIII have $N\approx10^3$ logical qubits and $D\approx10^9$--$10^{10}$ logical cycles, so they
need $\epsilon\approx10^{-13}$--$10^{-15}$ \emph{per logical qubit per cycle}. The best physical
two-qubit gates are at $10^{-3}$--$10^{-4}$. The twelve orders of magnitude in between are what
error correction has to buy, and what \S\ref{ph:sec:codes} is about.

\begin{keybox}[title={What a circuit asks of the physics}]
DiVincenzo's checklist, in the form the arithmetic needs it: (1) well-defined two-level systems
that can be made in quantity; (2) reset to $\ket0$; (3) single-qubit rotations about two axes
(a drive with a controllable phase); (4) one entangling two-qubit gate; (5) QND readout, fast
enough to be used mid-circuit; and (6) coherence times long compared with the gate time, since
the ratio $T_2/t_{\rm gate}$ bounds how many gates fit before the state is lost. The estimate of
\S\ref{rsa:sec:phys} then consumes three derived numbers: the error per operation $p$, the
duration of one round of parity checks (the \emph{cycle}), and the time from a measurement to
the next conditional operation (the \emph{reaction time}).
\end{keybox}

% ===========================================================================
\section{How the qubits are built}\label{ph:sec:hw}

Every platform below realises the same abstract object --- a two-level system with a Hamiltonian
we can switch --- in a different physical degree of freedom. The choice fixes the gate time, the
coherence time, the readout, the connectivity and the dominant noise, and through them the three
numbers of the estimate.

\subsection{Superconducting circuits}\label{ph:sec:sc}

\paragraph{Superconductivity.} Below a critical temperature $T_c$ (aluminium $1.2$\,K, niobium
$9.3$\,K, NbTi about $10$\,K) some metals lose all resistance. The BCS theory of 1957 explains
it: the exchange of lattice vibrations (phonons) gives electrons near the Fermi surface a weak
attraction that overcomes their Coulomb repulsion and binds them in \emph{Cooper pairs} of
opposite momentum and spin. A pair is a boson, so the pairs condense into a single macroscopic
quantum state with a well-defined phase $\varphi$. Breaking a pair costs an energy $2\Delta$ with
the gap $\Delta\approx1.76\,k_BT_c$; for aluminium $\Delta/h\approx44$\,GHz, so a $5$\,GHz photon
cannot break a pair and the condensate is protected from the qubit's own control signals. The
condensate expels magnetic fields (the Meissner effect) except within the London penetration depth
of $20$--$100$\,nm; type II superconductors such as NbTi tolerate strong fields by letting flux in
as quantised vortices, which is why NbTi is used for the readout cabling and MRI magnets, while the
qubits themselves use aluminium and niobium, whose isotropic ($s$-wave) pairing gives clean,
reproducible junctions \cite{ph:uqt26}.

\paragraph{The Josephson junction.} Two superconductors separated by a $\sim1$\,nm oxide barrier
exchange Cooper pairs by tunnelling. Josephson's relations (1962) are
\[ I=I_c\sin\varphi,\qquad V=\frac{\hbar}{2e}\frac{d\varphi}{dt}, \]
with $\varphi$ the phase difference across the barrier and $I_c$ the critical current. Combining
them, the junction behaves as a nonlinear inductor,
$L_J(\varphi)=\Phi_0/(2\pi I_c\cos\varphi)$ with $\Phi_0=h/2e$ the flux quantum, storing energy
$-E_J\cos\varphi$ with $E_J=\Phi_0I_c/2\pi$. This single nonlinear element is the whole trick: a
plain $LC$ circuit is a harmonic oscillator whose levels are equally spaced, so a drive at the
$0\to1$ frequency also drives $1\to2$ and nothing singles out two levels; the cosine potential
makes the spacing uneven and turns the circuit into an \emph{artificial atom}.

\paragraph{The transmon.} A junction shunted by a large capacitor $C$ has the Hamiltonian
\[ H=4E_C\,(\hat n-n_g)^2-E_J\cos\hat\varphi,\qquad E_C=\frac{e^2}{2C}, \]
where $\hat n$ counts Cooper pairs on the island and $[\hat\varphi,\hat n]=i$, the same
commutator as $x$ and $p$. In the transmon regime $E_J/E_C\approx50$--$100$ the phase is
well localised near $0$, so expand the cosine:
$-E_J\cos\hat\varphi=-E_J+\tfrac12E_J\hat\varphi^2-\tfrac1{24}E_J\hat\varphi^4+\dots$ The
quadratic term with $4E_C\hat n^2$ is a harmonic oscillator of plasma frequency
$\hbar\omega_p=\sqrt{8E_JE_C}$ (the $LC$ frequency with $L_J=\Phi_0^2/4\pi^2E_J$), and the
quartic term, treated in first-order perturbation theory with
$\hat\varphi=(2E_C/E_J)^{1/4}(a+a^\dagger)$, lowers each level $m$ by
$\tfrac{E_C}2(m^2+m)$. Hence
\[ \hbar\omega_{01}\approx\sqrt{8E_JE_C}-E_C,\qquad
   \alpha\equiv\omega_{12}-\omega_{01}\approx-E_C/\hbar , \]
and in the qubit subspace $H\approx-\tfrac{\hbar\omega_{01}}2Z$: the abstract qubit of
\S\ref{ph:sec:unitary}, with the drive term $\hbar\Omega\cos(\omega t)X$ supplied by a microwave
voltage capacitively coupled to the island, which couples to $\hat n\propto(a-a^\dagger)$, i.e.\ to
$Y$ or $X$ depending on its phase.
Typical values: $E_J/h\approx15$\,GHz, $E_C/h\approx250$\,MHz, so $\omega_{01}/2\pi\approx5$\,GHz
and the anharmonicity $\alpha/2\pi\approx-250$\,MHz. The large $E_J/E_C$ exponentially suppresses
the sensitivity to the offset charge $n_g$ (charge noise), which is what made the transmon
(2007) usable where the earlier charge qubit (1999) was not. The price is the small anharmonicity:
a $20$\,ns pulse has a $22$\,MHz bandwidth, an order of magnitude below $|\alpha|$, so leakage
to $\ket2$ is suppressed but not absent, and pulse shaping (DRAG, which adds a derivative
component to the envelope) brings it to the $10^{-4}$ level \cite{ph:uqt26}. Fluxonium replaces
the shunt capacitor by a large inductance, gains a much larger anharmonicity and $T_1>1$\,ms, at
the cost of a sub-GHz qubit frequency and more complex control; Google acquired a fluxonium
company in 2025.

\paragraph{Why 5\,GHz and why 15\,mK.} A qubit at $5$\,GHz has $hf/k_B\approx0.24$\,K. At the
$15$\,mK of a dilution refrigerator the thermal population of $\ket1$ is $e^{-hf/k_BT}\approx
e^{-16}\approx10^{-7}$, so the qubit initialises itself by waiting; at $100$\,mK it would be
$10^{-5}$ and at $1$\,K, $9\%$. Lower frequencies would need colder stages; higher ones run into
dielectric and quasiparticle losses and the $4$--$8$\,GHz band of available microwave components.
The refrigerator is a $^3$He--$^4$He dilution unit: below $0.87$\,K the mixture separates into a
$^3$He-rich phase and a dilute phase, and $^3$He crossing the boundary absorbs heat, a process
that keeps working to the millikelvin range because the solubility of $^3$He in superfluid $^4$He
stays finite ($6.6\%$) as $T\to0$. The cooling power at the mixing chamber scales as
$\dot Q\approx84\,\dot n_3T^2$ and is only $10$--$30\,\mu$W at $20$\,mK, against $1$--$2$\,W at the
$4$\,K stage; a single $15$\,mK cryostat draws about $16$\,kW from the wall. Every microwave line
into it is attenuated by $60$\,dB in stages so that only a few hundred thermal photons' worth of
noise reaches the chip, and every readout line leaves through a quantum-limited parametric
amplifier (JPA or TWPA) at the base stage, a HEMT at $4$\,K and a room-temperature amplifier
\cite{ph:uqt26}. This wiring, one to two lines per qubit, is the scaling constraint that
multiplexing and cryo-CMOS are trying to lift.

\paragraph{Gates.} Single-qubit gates are the Rabi drive of \S\ref{ph:sec:rabi}: an arbitrary
waveform generator shapes an envelope, a mixer puts it on a $4$--$8$\,GHz local oscillator with
in-phase and quadrature components that set the rotation axis, and the pulse lasts $10$--$40$\,ns.
Two-qubit gates use a coupling element between neighbours. With a \emph{tunable coupler} (Google
since Sycamore, IBM since Heron in 2023) a flux pulse briefly switches on a $ZZ$ interaction
or brings the qubits into resonance for an iSWAP; the gate lasts $20$--$200$\,ns. The older
cross-resonance gate drives the control qubit at the target's frequency and lasts longer. Readout
is \emph{dispersive}: a readout resonator at $\omega_r$ coupled to the qubit with strength $g$ at
detuning $\Delta=\omega_{01}-\omega_r$ has its frequency shifted by $\pm\chi$ with
$\chi\approx g^2/\Delta$ depending on the qubit state (the Jaynes--Cummings Hamiltonian
$H=\hbar\omega_ra^\dagger a+\tfrac{\hbar\omega_{01}}2Z+\hbar g(a^\dagger\sigma_-+a\sigma_+)$ in the
dispersive limit $g\ll\Delta$). A microwave tone reflected off the resonator picks up a
state-dependent phase, and the measurement is QND because it only touches $Z$. It takes
$50$--$500$\,ns including the resonator ring-down; a Purcell filter stops the qubit from decaying
through the readout line \cite{ph:uqt26}.

\paragraph{Noise.} Energy relaxation is dominated by two-level-system (TLS) defects in the
amorphous oxides and interfaces, by quasiparticles, and by the Purcell channel; dephasing by $1/f$
flux and charge noise; and the hardest problem for a code is \emph{correlated} errors: a cosmic
ray or a radioactive decay in the substrate creates a burst of phonons that breaks Cooper pairs
across a millimetre of chip and flips many qubits at once, a few times a minute. Gap engineering
(making the junction electrodes' gaps unequal so quasiparticles are trapped away from the
junction), phonon traps and shielding have brought the rate down but not to zero
\cite{ph:uqt26}.

\paragraph{Where it stood in 2026.} $T_1$ of $100$--$400\,\mu$s for transmons, above $1$\,ms for
fluxonium and $1.67$\,ms for a tantalum-on-silicon transmon at Princeton; single-qubit errors
$\approx10^{-4}$; two-qubit errors: $99.85\%$ median CZ on IBM's 156-qubit Heron r3 and $99.67\%$
on Google's 105-qubit Willow, with the best pairs below $10^{-3}$; readout $99.4\%$ in $2.4\,\mu$s
on Heron-class devices; chips of $105$--$156$ qubits with four nearest neighbours (Google since
2019, IBM's 120-qubit Nighthawk since 2025) and IBM's Loon test chip adding two long-range
couplers per qubit for the qLDPC codes of \S\ref{ph:sec:codes}. The surface-code cycle ---
four layers of CZ gates, a readout and a reset --- is about $1\,\mu$s on Willow, which is the
cycle time of \S\ref{rsa:sec:phys} \cite{ph:uqt26,ph:willow24}.

\subsection{Trapped ions}\label{ph:sec:ions}

\paragraph{The trap.} An ion is a charged atom, so an electric field holds it. A static field
cannot confine a charge in all three directions (Earnshaw's theorem), so a Paul trap uses a
radio-frequency quadrupole field at $\Omega_{\rm rf}/2\pi\approx20$--$50$\,MHz whose time average
is a harmonic pseudo-potential with secular frequencies $\nu/2\pi\approx1$--$5$\,MHz. Modern traps
are microfabricated surface-electrode chips with the ions $50$--$100\,\mu$m above the surface,
segmented so that DC voltages can split, move and merge ion chains between zones: the QCCD
architecture used by Quantinuum, Oxford Ionics (IonQ since 2025) and Universal Quantum. Ions in a
chain repel each other and sit $2$--$5\,\mu$m apart; their motion is a set of collective normal
modes, each a quantum harmonic oscillator with phonon number $n$, and that shared motion is the
bus for two-qubit gates \cite{ph:uqt26}.

\paragraph{The qubit.} Two long-lived internal states of the ion. Hyperfine qubits use the
$F=0$ and $F=1$, $m_F=0$ ground states of an odd isotope with nuclear spin, split by
$12.64$\,GHz in $^{171}$Yb$^+$ and $3.2$\,GHz in $^{43}$Ca$^+$; at a magnetic field where
$d\nu/dB=0$ (a \emph{clock} transition) they are first-order insensitive to field noise and
coherence times of $T_2\approx70$\,s have been measured, a coherence-to-gate ratio above $10^6$.
Optical qubits (AQT) use a metastable $D$ level $\approx400$\,THz above the ground state.
Because ions of one isotope are identical, there is no fabrication spread of qubit frequencies;
the variability sits in the control fields.

\paragraph{Cooling and the Lamb--Dicke regime.} Gates need the motion cold. Doppler cooling
illuminates the ion with light tuned just below a cycling transition of linewidth $\Gamma$: an ion
moving towards the beam sees the light Doppler-shifted into resonance, absorbs a photon and its
momentum $\hbar k$ against its motion, and re-emits in a random direction, so on average it slows.
The limit $T_D=\hbar\Gamma/2k_B$ is about $0.5$\,mK for Yb$^+$ ($145\,\mu$K for the Rb atoms of
\S\ref{ph:sec:atoms}). Resolved-sideband cooling then removes phonons one at a time by driving the
$\ket{g,n}\to\ket{e,n-1}$ transition. The gates work in the Lamb--Dicke regime
$\eta\sqrt{2\bar n+1}\ll1$, $\eta=k\sqrt{\hbar/2m\nu}\approx0.05$--$0.1$, where the ion's
excursion is small compared with the laser wavelength and carrier and sideband transitions are
cleanly resolved.

\paragraph{Gates.} Single-qubit rotations are microwave or Raman $\pi$ pulses of $1$--$50\,\mu$s.
The M\o lmer--S\o rensen gate is worth deriving, because it is the cleanest example of an
entangling Hamiltonian that is not simply ``switched on''. A laser on ion $j$ couples its spin to
the shared mode through the Lamb--Dicke factor: in the interaction picture
$H_j=\hbar\Omega\,\sigma_x^{(j)}\cos\big(\omega_Lt-\eta(ae^{-i\nu t}+a^\dagger e^{i\nu t})\big)$,
and to first order in $\eta$ a bichromatic drive at $\omega_0\pm(\nu+\delta)$ leaves, after the
rotating-wave approximation,
\[ H_{\rm MS}=\hbar\eta\Omega\,\big(\sigma_x^{(1)}+\sigma_x^{(2)}\big)\big(a\,e^{i\delta t}+a^\dagger e^{-i\delta t}\big) . \]
This is a force on the oscillator whose sign depends on the collective spin $S_x=\sigma_x^{(1)}+\sigma_x^{(2)}$.
Its exact propagator is a displacement $D(\alpha(t)S_x)$ with
$\alpha(t)=(\eta\Omega/\delta)(1-e^{-i\delta t})$, a circle in phase space, times a phase
$e^{i\Phi(t)S_x^2}$ with $\Phi(t)=(\eta\Omega/\delta)^2(\delta t-\sin\delta t)$. At
$t=2\pi/\delta$ the circle closes, the displacement vanishes whatever the initial phonon number,
and what remains is $e^{i\Phi S_x^2}$ with $\Phi=2\pi(\eta\Omega/\delta)^2$; since
$S_x^2=2I+2\sigma_x^{(1)}\sigma_x^{(2)}$, choosing $\eta\Omega/\delta=1/\sqrt8$ gives
$\Phi=\pi/4$ and the gate $e^{i\frac\pi4X\otimes X}$ of \S\ref{ph:sec:unitary}, up to a global
phase. The entanglement is a \emph{geometric} phase, the area enclosed by the loop; the motion is
only borrowed. Residual spin--motion entanglement from an unclosed loop, mode-frequency drift
and laser phase noise are the gate's error sources. Any pair in a chain can be gated this way --- effective all-to-all connectivity --- but
the gate takes $10$--$100\,\mu$s (longer for larger chains, up to a millisecond on IonQ's Forte),
and shuttling between zones adds milliseconds \cite{ph:uqt26}. Readout is state-dependent
fluorescence: one qubit state scatters photons on a cycling transition, the other is dark;
collecting enough photons takes $100\,\mu$s to a few ms on cameras, $11\,\mu$s at $99.93\%$ with a
trap-integrated superconducting nanowire detector.

\paragraph{Where it stood in 2026.} The best fidelities of any platform: single-qubit error
$1.5\times10^{-7}$ on a $^{43}$Ca$^+$ clock qubit (Oxford--Osaka, 2025), two-qubit $99.92\%$
averaged over all pairs of Quantinuum's 98-qubit Helios (November 2025), $>99.99\%$ on a two-ion
Oxford Ionics chip, SPAM error $9\times10^{-5}$ in $^{137}$Ba$^+$. Helios is the largest
gate-based ion processor; $512$-ion crystals are controlled for analog simulation
\cite{ph:uqt26,ph:helios}. The limitation is the clock: with gates of tens of microseconds and
readout of hundreds, one round of parity checks takes of the order of a millisecond rather than a
microsecond, so the same logical algorithm runs a thousand times longer unless connectivity and
transversal gates recover part of it (\S\ref{ph:sec:feeds}).

\subsection{Neutral atoms}\label{ph:sec:atoms}

\paragraph{Traps and arrays.} A neutral atom has no charge to hold, so it is held by light. A
far-detuned laser induces a dipole, and the AC Stark shift $U=-\tfrac12\alpha|E|^2$ with
polarisability $\alpha>0$ for red detuning pulls the atom to the focus: an \emph{optical
tweezer}. A spatial light modulator imprints a hologram that makes hundreds to thousands of
tweezers at once, a magneto-optical trap loads them from a vapour at random (about $50\%$
occupancy), and acousto-optic deflectors move atoms to fill a defect-free array, in $2$D or $3$D,
$3$--$5\,\mu$m apart. The atoms are laser-cooled to a few $\mu$K; the chamber can be at room
temperature, and the record arrays --- $6,100$ caesium atoms in $12,000$ tweezers at Caltech in
2024, with $12.6$\,s hyperfine coherence and $99.98\%$ single-qubit gates --- were
\cite{ph:uqt26,ph:caltech6100}.

\paragraph{The qubit and the Rydberg blockade.} The qubit is a pair of hyperfine ground states
($^{87}$Rb, $^{133}$Cs) or nuclear-spin states ($^{171}$Yb, $^{87}$Sr), coherent for seconds.
The interaction is borrowed temporarily: a laser excites an atom to a \emph{Rydberg state} of
principal quantum number $n\approx50$--$100$, whose orbital radius grows as $n^2a_0$ (about
$0.5\,\mu$m at $n=100$), whose lifetime grows as $n^3$ (tens to hundreds of $\mu$s), and whose
van der Waals interaction with another Rydberg atom at distance $r$ is $V=C_6/r^6$ with
$C_6\propto n^{11}$. When $V$ exceeds the laser's Rabi frequency $\hbar\Omega$, two atoms within
the blockade radius
\[ R_b=\big(C_6/\hbar\Omega\big)^{1/6}\approx5\text{--}10\,\mu\text{m} \]
cannot both be excited: the doubly excited state is shifted out of resonance. In the basis $\{\ket0,\ket1,\ket r\}$ per atom, with the laser coupling only
$\ket1\leftrightarrow\ket r$,
\[ H=\frac{\hbar\Omega}2\sum_{j=1,2}\big(\ket1\bra r_j+\ket r\bra1_j\big)+V\ket{rr}\bra{rr},\qquad V\gg\hbar\Omega . \]
Start in $\ket{00}$: nothing couples, no phase. In $\ket{01}$ or $\ket{10}$ one atom Rabi-cycles
$\ket1\to\ket r\to\ket1$ in a $2\pi$ pulse and picks up the phase $-1$ of a full spin rotation.
In $\ket{11}$ the laser couples $\ket{11}$ only to the symmetric state
$(\ket{1r}+\ket{r1})/\sqrt2$ with enhanced Rabi frequency $\sqrt2\Omega$, because $\ket{rr}$ is
blocked by $V$; the Levine--Pichler protocol (2019) uses two pulses with a phase jump between
them so that both the singly excited and the doubly blocked branches return home with phases
that differ by exactly $\pi$: the result is $\operatorname{diag}(1,e^{i\phi},e^{i\phi},e^{i(2\phi+\pi)})$,
a $CZ$ after single-qubit $Z$ rotations. The blockade is the $ZZ$ energy shift of
\S\ref{ph:sec:unitary} in disguise, and it needs no individual addressing: a global beam gates
every pair placed within $R_b$ in parallel, which is how 60 pairs were gated at once in 2023. The pulse lasts $0.2$--$2\,\mu$s; the fidelity reached $99.5\%$ on 60
atoms in parallel in 2023 and $99.85\%$ (raw) in a 2026 zoned architecture \cite{ph:uqt26,ph:evered23}.

\paragraph{Architecture.} The decisive advantage is that the atoms move. The Harvard--MIT--QuEra
experiment of December 2023 kept data atoms in a storage zone, carried them with tweezers into an
entangling zone, and read ancillas in a separate zone: $280$ atoms, up to $48$ logical qubits,
transversal logical CNOTs by moving whole code blocks next to each other
\cite{ph:bluvstein23}. Movement gives any connectivity at all, so codes that need long-range
checks (qLDPC, \S\ref{ph:sec:codes}) are natural here, and the cost is time: coherent transport,
cooling and camera imaging take milliseconds even though the gate pulse is sub-microsecond, so a
logical cycle is again of order a millisecond. Alkaline-earth atoms add a second gift: decay out of
the qubit manifold can be detected (\emph{erasure conversion}), and an error whose location is
known is far cheaper to correct than one whose location must be inferred. Atom loss is the other
specific error, replenished from a reservoir in the newer architectures.

\subsection{Spins in silicon}\label{ph:sec:spins}

An electron (or hole) confined in a gate-defined quantum dot, $\sim50$\,nm across, has a spin
$\tfrac12$ whose degeneracy a static field $B$ lifts by the Zeeman energy
$E_Z=g\mu_BB$: at $B=1$\,T with $g\approx2$, $E_Z=116\,\mu$eV, a $28$\,GHz qubit. This is
larger than $k_BT$ at $1$\,K ($86\,\mu$eV), so spin qubits can run at $0.5$--$1.5$\,K where the
cryostat has milliwatts of cooling power instead of microwatts, enough for control electronics next
to the qubits. Single-qubit gates are electron spin resonance (an oscillating magnetic field) or
electric-dipole spin resonance (an oscillating gate voltage converted to a magnetic drive by
spin--orbit coupling or a micromagnet gradient); two-qubit gates lower the tunnel barrier between
two dots so that the exchange interaction $J\vec S_1\cdot\vec S_2$ of \S\ref{ph:sec:unitary} acts
for $\sim75$\,ns --- $J$ is not a magnetic force but a consequence of the Pauli principle: with
tunnelling $t$ between the dots and charging energy $U$ for two electrons on one dot, second-order
perturbation theory of the Hubbard model gives $J\approx4t^2/U$, non-zero only for the spin
singlet, whose antisymmetric spin part allows the symmetric orbital part that tunnels, so a
voltage on the barrier gate tunes $t$ and switches $J$ on and off by orders of magnitude; readout converts spin to charge (Pauli spin blockade: a second electron can
tunnel onto the dot only if the two spins form a singlet) and senses the charge with a
single-electron transistor or gate reflectometry. Isotopically purified $^{28}$Si removes the
nuclear spins that dephase the electron, and the dominant noise is then charge noise in the oxide.
Fidelities of $99.5\%$ for two-qubit and $99.8$--$99.9\%$ for single-qubit gates have been reached
on two-qubit devices, and $>99\%$ on unit cells from a $300$\,mm foundry; the largest arrays are
$18$ qubits (HRL exchange-only, TU Delft, 2026). The appeal is density --- gate pitches below
$100$\,nm and CMOS manufacturing --- and the open problems are variability, wiring and
demonstrated entanglement beyond a handful of qubits \cite{ph:uqt26}.

\subsection{Photons}\label{ph:sec:photons}

A photon carries a qubit in its polarisation, its path, its arrival time bin, or its frequency,
and flies at $c$ through waveguides, so it decoheres only by being lost. Single-qubit gates are
free: a half-wave plate or a Mach--Zehnder interferometer with a phase shifter is an arbitrary
rotation. The difficulty is the opposite of the matter platforms': photons do not interact, so a
two-qubit gate cannot be a Hamiltonian switched on. Linear optics with measurement (KLM, 2001)
makes probabilistic entangling gates, and the practical architectures sidestep gates altogether:
\emph{measurement-based} computing prepares a large entangled cluster state and computes by
measuring it in chosen bases, and PsiQuantum's \emph{fusion-based} variant builds the cluster from
small resource states fused by Bell measurements, each of which succeeds with probability $1/2$
unless boosted. Everything then depends on components: single-photon sources (quantum dots at
$4$\,K with $>99\%$ indistinguishability and $57$--$70\%$ end-to-end efficiency; parametric
down-conversion pairs at room temperature but probabilistic) and detectors (superconducting
nanowires at $1$--$3$\,K, $>95\%$ efficiency, sub-nanosecond timing). A photon lost is an
\emph{erasure}, located by the detector that did not fire, and the fault-tolerance thresholds
for loss are a few percent per component. The continuous-variable branch (Xanadu) encodes
qubits in field quadratures as GKP states and needs $>10$\,dB of squeezing, against a record of
$15$\,dB measured. The largest demonstrations in 2026 were $12$-mode processors and
$99.22\%$ fusions on PsiQuantum's test chips; the largest fault-tolerant plans are helium-cooled
data-centre buildings, not chandeliers \cite{ph:uqt26}.

\subsection{Three others}

\paragraph{Bosonic cat qubits.} A microwave cavity holds a superposition of two coherent states
$\ket{\pm\alpha}$ stabilised by two-photon dissipation. Bit flips between $\ket\alpha$ and
$\ket{-\alpha}$ are exponentially suppressed in the photon number ($T_1$ for bit flips of
$22$\,s reported), while phase flips grow only linearly, so the noise is \emph{biased} by
$\eta=p_Z/p_X>10^5$. A one-dimensional repetition code then suffices for the phase flips, and
the Alice \& Bob group's estimates are $349,133$ cat qubits for RSA-2048 in four days (2019) and
$126,133$ for a 256-bit elliptic curve in nine hours (2023) \cite{ph:cats19,ph:cats23}.

\paragraph{NV centres.} A nitrogen-vacancy defect in diamond is a spin-1 electron system
readable optically at room temperature, with nearby $^{13}$C and $^{14}$N nuclear spins as
long-lived register qubits. It is the platform of the 2015 loophole-free Bell test and of quantum
sensing; as a processor it has stayed below ten qubits for years.

\paragraph{Majorana (topological) qubits.} Zero-energy modes at the ends of a one-dimensional
topological superconductor would store a qubit non-locally in the fermion parity of a pair, and
braiding them would implement Clifford gates with topological protection. Microsoft's 2025--2026
devices measure parities in gate-defined InAs--Al nanowires, but as of September 2026 no operational
qubit metrics (coherence, gate fidelities) had been published, earlier Majorana claims had been
retracted, and the $T$ gate is not protected even in principle, so a code is still needed
\cite{ph:uqt26}.

\subsection{The platforms side by side, and what they feed into the estimate}\label{ph:sec:feeds}

\begin{center}\footnotesize
\renewcommand{\arraystretch}{1.15}
\begin{tabular}{@{}p{2.4cm}p{2.5cm}p{2.5cm}p{2.5cm}p{2.5cm}p{2.3cm}@{}}
\toprule
 & superconducting & trapped ions & neutral atoms & silicon spins & photons\\
\midrule
qubit & two levels of a nonlinear $LC$ circuit & hyperfine or optical states of an ion & hyperfine or nuclear states of an atom & electron/hole spin in a quantum dot & polarisation, path or time bin\\
entangling gate & tunable coupler $ZZ$/iSWAP & M\o lmer--S\o rensen via phonons & Rydberg blockade $CZ$ & exchange $\sqrt{\rm SWAP}$/CZ & fusion (probabilistic)\\
2-qubit time & $20$--$200$\,ns & $10$--$1000\,\mu$s & $0.2$--$2\,\mu$s (+ms transport) & $40$--$75$\,ns & ns (circuit traversal)\\
best 2-qubit fidelity, commercial & $99.85\%$ (IBM Heron r3), $99.67\%$ (Willow) & $99.92\%$ (Helios, 98 q.) & $99.5$--$99.85\%$ & $99.5\%$ (2 q.) & $99.22\%$ fusion\\
$T_1$ / coherence & $0.1$--$1.7$\,ms & seconds to hours & $10$\,s--$1$\,h trap; $12.6$\,s $T_2$ & $>1$\,s; ms $T_2$ & loss-limited\\
readout & $0.1$--$2.4\,\mu$s dispersive & $11\,\mu$s--ms fluorescence & ms imaging ($2.4\,\mu$s record) & $0.2$--few $\mu$s reflectometry & $<1$\,ns detector\\
operating $T$ & $15$\,mK & ions $<1$\,mK, trap RT--$4$\,K & atoms $\mu$K, chamber RT & $0.1$--$1.5$\,K & RT circuits, $2$--$4$\,K sources/detectors\\
connectivity & 4 neighbours (+2 long range on Loon) & all-to-all in a zone & any, by moving atoms & neighbours (+shuttling) & any, by routing\\
largest system, 2026 & 156 (IBM), 105 (Google) & 98 (Quantinuum) & 256 gated; $6,100$ trapped & 18 & 12 modes\\
\bottomrule
\end{tabular}
\end{center}
\noindent Figures as reported by vendors and papers to September 2026, compiled in
\cite{ph:uqt26} (Tables 17--18, 26); laboratory records on two qubits are better in every column.

The estimate of \S\ref{rsa:sec:phys} reads this table through three numbers:
\begin{itemize}\setlength\itemsep{2pt}
\item \textbf{$p$} is the error per \emph{location} of the syndrome circuit --- gates, idles,
  measurements and resets alike --- at scale, with crosstalk, not the best pair in a lab. The
  model's $p=10^{-3}$ sits at the median two-qubit error of the 2026 superconducting chips and
  somewhat above the ion and atom figures. Halving $p$ to $5\times10^{-4}$ drops the distance from
  $25$ to $21$ and the qubits per logical qubit from $1,352$ to $968$ (\S\ref{rsa:sec:phys}); at
  $10^{-4}$ it is $392$. The sensitivity is exponential in the exponent and only polynomial in
  the prefactor.
\item \textbf{The cycle} is one round of parity checks: four entangling layers, a readout and a
  reset. Superconducting chips do it in $\approx1\,\mu$s; ions and atoms, with microsecond gates
  but $10^2$--$10^3\,\mu$s readouts and transport, in about a millisecond. Runtime is linear in
  the cycle, so the $15$ minutes for ECDLP-256 and $5$ days for RSA-2048 become ten days and
  thirteen years at a $1$\,ms cycle, unless the architecture compensates. That is why the slow
  platforms lean on all-to-all connectivity, transversal gates and higher-rate codes: IonQ's
  2026 estimate of $26$ days for a 256-bit curve on $20,000$ ions \cite{ec:ionq26} buys its time
  back with those, and still reports a $63\%$ success probability.
\item \textbf{The reaction time} is decoder latency plus control: the $10\,\mu$s of the model is
  the time between an ancilla measurement and the next operation that depends on it, which
  matters because every $T$ gate ends in a measurement whose outcome decides a Clifford
  correction (\S\ref{ph:sec:ft}). It is the only one of the three numbers set by classical
  electronics rather than physics.
\end{itemize}
Connectivity enters through the code family and the routing overhead rather than through a
number, and leakage, loss and correlated errors enter through whatever the chosen code does not
model: the estimates assume independent Pauli errors, and the hardware does not quite deliver
them.

\subsection{The challenges, platform by platform}

What stands between each platform and the thousand $d=25$ patches of the estimate, in the
order the field itself lists them \cite{ph:uqt26} (Table 17):
\begin{itemize}\setlength\itemsep{2pt}
\item \textbf{Superconducting.} Fidelity uniformity across a chip (the median, not the best
  pair, sets $p$); crosstalk and frequency crowding growing with qubit count; correlated errors
  from radiation; leakage; the wiring, cooling budget and energy of $10^6$ control lines ---
  hence cryo-CMOS, superconducting (SFQ) logic, multiplexed readout and chip-to-chip microwave
  links; and two-qubit errors improving more slowly than single-qubit ones. The breakthroughs
  were the transmon (2007), tunable couplers (2019), the Willow below-threshold memory (2024) and
  gap engineering against quasiparticles (2024).
\item \textbf{Trapped ions.} Gate and readout speed; control of many motional modes as chains
  grow; anomalous heating of the motion from the trap surface; laser power, phase noise and
  beam-pointing stability; shuttling overhead. The breakthroughs were the QCCD architecture,
  microwave-driven gates (no lasers), $10^{-7}$ single-qubit errors (2025) and the 98-ion Helios
  with all-pairs $99.92\%$ (2025).
\item \textbf{Neutral atoms.} Loading, rearrangement and imaging times that dominate the duty
  cycle; atom loss and leakage; optical power and uniformity across thousands of tweezers;
  QND mid-circuit readout without disturbing neighbours. The breakthroughs were parallel
  high-fidelity Rydberg gates (2023), zoned logical processing (2023), erasure conversion in
  ytterbium (2023--2025) and $6,100$-atom arrays (2024).
\item \textbf{Silicon spins.} Charge noise and device-to-device variability, which forces
  per-qubit calibration; valley splitting in silicon; demonstrated entanglement only on a
  handful of dots; fabrication cycles of months. The breakthroughs were $>99\%$ two-qubit gates
  (2022), operation above $1$\,K (2020) and foundry-made $300$\,mm unit cells (2025).
\item \textbf{Photons.} Source efficiency and indistinguishability; loss in every component;
  probabilistic entanglement that multiplies the resource states needed; large cluster states;
  cryogenic detector integration. The breakthroughs were bright quantum-dot sources, fusion-based
  architectures and $>99\%$ fusion on test chips (2025--2026).
\end{itemize}
Common to all: no platform had, by September 2026, run a logical circuit with real-time decoding
and without post-selection at a distance above $7$, and every resource estimate assumes a noise
model (independent, Pauli, uniform at $10^{-3}$) that the hardware does not yet deliver.

% ===========================================================================
\section{Errors, codes, and the price of a logical qubit}\label{ph:sec:codes}

\S\ref{rsa:sec:phys} used the surface code as a black box with one formula. This section opens
it: how parity checks locate errors without touching the data, why the logical error falls
exponentially with the distance below a threshold, what the decoder must do in real time, and
how logical gates --- above all the $T$ gate --- are performed. The last subsection returns the
favour by deriving the ``one CCZ every $25\,\mu$s'' that drives every runtime in Part~VIII.

\subsection{Stabilisers: measuring parity, not data}

A stabiliser code on $n$ qubits is defined by $n-k$ commuting Pauli operators $S_j$; the code
space is their joint $+1$ eigenspace and holds $k$ logical qubits. Each $S_j$ is measured by an
ancilla: for $S=Z_1Z_2$, two CNOTs from the data qubits onto an ancilla in $\ket0$ and a readout
give the parity of the two bits without revealing either bit, since $\ket{00}$ and $\ket{11}$ both
leave the ancilla at $0$. A Pauli error $E$ that anticommutes with $S_j$ flips its outcome; the
pattern of flipped outcomes is the \emph{syndrome}. Because the measurement projects
(\S\ref{ph:sec:open}), a continuum of errors becomes a Pauli with a syndrome.

The rotated surface code of distance $d$ places $d^2$ data qubits on a grid with $d^2-1$ ancillas
in the faces: each face measures the parity of its four corners, in the $Z$ basis on half the
faces (catching $X$ errors) and the $X$ basis on the other half (catching $Z$ errors), with
weight-two checks along the edges. It is $[[d^2,1,d]]$ with $2d^2-1$ physical qubits.

\begin{center}
\begin{tikzpicture}[x=11mm,y=11mm]
  \fill[cblue!14,draw=cblue!55] (0,-1) rectangle (1,0);
  \fill[cblue!14,draw=cblue!55] (1,-2) rectangle (2,-1);
  \fill[cverm!12,draw=cverm!50] (1,-1) rectangle (2,0);
  \fill[cverm!12,draw=cverm!50] (0,-2) rectangle (1,-1);
  \fill[cverm!12,draw=cverm!50] (0,0) arc (180:0:0.5) -- cycle;
  \fill[cverm!12,draw=cverm!50] (1,-2) arc (180:360:0.5) -- cycle;
  \fill[cblue!14,draw=cblue!55] (0,-1) arc (90:270:0.5) -- cycle;
  \fill[cblue!14,draw=cblue!55] (2,0) arc (90:-90:0.5) -- cycle;
  \node[font=\scriptsize,color=cblue] at (0.5,-0.5) {$z_1$};
  \node[font=\scriptsize,color=cblue] at (1.5,-1.5) {$z_2$};
  \node[font=\scriptsize,color=cblue] at (-0.3,-1.5) {$z_3$};
  \node[font=\scriptsize,color=cblue] at (2.3,-0.5) {$z_4$};
  \foreach \c/\r/\q in {0/0/1,1/0/2,2/0/3,0/1/4,1/1/5,2/1/6,0/2/7,1/2/8,2/2/9}
    \node[circle,fill=white,draw=black,inner sep=0.8pt,minimum size=4.2mm,font=\scriptsize] at (\c,-\r) {\q};
\end{tikzpicture}
\hspace{1.2cm}
\begin{tikzpicture}[x=7.6mm,y=7.6mm,baseline=(current bounding box.center)]
  \foreach \i/\j in {0/0,0/2,1/1,1/3,2/0,2/2,3/1,3/3} \fill[cblue!12,draw=cblue!45] (\i,\j) rectangle ++(1,1);
  \foreach \i/\j in {0/1,0/3,1/0,1/2,2/1,2/3,3/0,3/2} \fill[cverm!12,draw=cverm!45] (\i,\j) rectangle ++(1,1);
  \foreach \i/\j in {1/1,2/2} \draw[cblue,line width=1.6pt,fill=cblue!35] (\i,\j) rectangle ++(1,1);
  \draw[cpurple,line width=2.2pt] (-0.15,0) -- (4.15,0);
  \foreach \i in {0,...,4} \foreach \j in {0,...,4} \fill[black] (\i,\j) circle (1.6pt);
  \node[draw=cverm,fill=white,circle,inner sep=0.6pt,font=\scriptsize\bfseries,text=cverm] at (2,2) {$X$};
  \node[font=\scriptsize,color=cpurple,anchor=north] at (2,-0.12) {logical $\bar X$};
\end{tikzpicture}
\end{center}
{\small\noindent Left: $d=3$, nine data qubits; \textcolor{cblue}{blue} $Z$-checks catch $X$
errors, \textcolor{cverm}{red} $X$-checks catch $Z$ errors. Right: $d=5$ (data qubits at the
vertices, checks on the faces): one $X$ error lights the two $Z$-faces beside it; the logical
$\bar X$ is a string of $d$ Paulis across the patch.}

\medskip
Written as a parity-check matrix for the $d=3$ patch, the $Z$-checks are
\[ H_Z=\begin{array}{r|ccccccccc}
    & 1&2&3&4&5&6&7&8&9\\ \hline
  z_1 & 1&1&0&1&1&0&0&0&0\\
  z_2 & 0&0&0&0&1&1&0&1&1\\
  z_3 & 0&0&0&1&0&0&1&0&0\\
  z_4 & 0&0&1&0&0&1&0&0&0\end{array}\qquad s=H_Z\,e\pmod2 . \]
An $X$ error on qubit 5 lights $z_1$ and $z_2$, and qubit 5 is the only one in both; an error on
qubit 7 lights $z_3$ alone. Errors on qubits 1 and 2 give the same syndrome, but $X_1X_2$
commutes with every check and acts trivially on the code space --- the two corrections are
equivalent modulo a stabiliser (\emph{degeneracy}), so either works. An undetectable error must
commute with all checks and act non-trivially: the shortest is a string of $d$ Paulis across the
patch, which is the logical operator itself. Hence the distance $d$, and the code corrects any
$\lfloor(d-1)/2\rfloor$ errors.

\subsection{The stabiliser formalism, with the mathematics}\label{ph:sec:stabmath}

\paragraph{The group.} The Pauli group on $n$ qubits, $\mathcal P_n$, is the set of tensor
products of $I,X,Y,Z$ with phases $\pm1,\pm i$. Two of its elements either commute or
anticommute, and the commutation is computed by a symplectic inner product: write
$P=X^{\mathbf a}Z^{\mathbf b}$ with bit vectors $\mathbf a,\mathbf b\in\mathbb F_2^n$; then
$P$ and $P'=X^{\mathbf a'}Z^{\mathbf b'}$ commute iff
$\mathbf a\cdot\mathbf b'+\mathbf b\cdot\mathbf a'=0\pmod2$. This turns every question about a
code into linear algebra over $\mathbb F_2$, which is why the parity-check matrix $H_Z$ above
told the whole story.

\paragraph{The code.} A stabiliser group $\mathcal S=\langle S_1,\dots,S_{n-k}\rangle$ is an
abelian subgroup of $\mathcal P_n$ not containing $-I$. The code space
$\mathcal C=\{\ket\psi:S\ket\psi=\ket\psi\ \forall S\in\mathcal S\}$ has dimension $2^k$: each
independent generator halves the space. The logical operators are the Paulis that commute with
all of $\mathcal S$ but are not in it, the \emph{normaliser} $N(\mathcal S)\setminus\mathcal S$;
they come in $k$ anticommuting pairs $(\bar X_i,\bar Z_i)$, and the distance $d$ is the minimum
weight among them. For the $d=3$ patch of the figure, $\bar Z=Z_1Z_2Z_3$ (a row) and
$\bar X=X_1X_4X_7$ (a column) commute with every check, anticommute with each other, and have
weight $3$; multiplying $\bar Z$ by the check $z_1=Z_1Z_2Z_4Z_5$ gives $Z_3Z_4Z_5$, another
representative of the same logical operator, which is the degeneracy seen earlier from the other
side.

\paragraph{Why errors are detected without disturbing the data.} Let $E\in\mathcal P_n$. For
$\ket\psi\in\mathcal C$, $S_j(E\ket\psi)=\pm E S_j\ket\psi=\pm E\ket\psi$, with the sign
$-1$ iff $E$ anticommutes with $S_j$. So $E\ket\psi$ is still a joint eigenvector of all the
checks, with a sign pattern (the syndrome) that depends on $E$ and not on $\ket\psi$. Measuring
the checks therefore learns $E$'s syndrome and nothing about the encoded state: the data is
untouched because the code space is a common eigenspace and the measurement commutes with every
logical operator. This is the precise content of ``measure parities, never the data''.

\paragraph{Knill--Laflamme.} A set of errors $\{E_a\}$ is correctable iff, with $P$ the
projector onto $\mathcal C$,
\[ P\,E_a^\dagger E_b\,P=c_{ab}\,P\qquad\text{for all }a,b , \]
for some matrix of scalars $c_{ab}$: no error product may map one codeword onto a different
codeword, nor distinguish codewords. For a stabiliser code this holds iff every product
$E_a^\dagger E_b$ is either in $\mathcal S$ (then $c_{ab}=\pm1$, the errors are degenerate) or
outside the normaliser (then it anticommutes with some check and $c_{ab}=0$). All Paulis of weight
$\le t$ satisfy it iff $2t<d$, and by linearity so does every operator in their span, which is
how the continuum of \S\ref{ph:sec:open} is covered.

\paragraph{The extraction circuit.} One weight-four $Z$-check and one $X$-check of the surface
code, as actually run every cycle:
\begin{center}
\begin{quantikz}[row sep=0.18cm,column sep=0.3cm]
   \lstick{$d_1$} & \ctrl{4} & \qw & \qw & \qw & \qw \\
   \lstick{$d_2$} & \qw & \ctrl{3} & \qw & \qw & \qw \\
   \lstick{$d_3$} & \qw & \qw & \ctrl{2} & \qw & \qw \\
   \lstick{$d_4$} & \qw & \qw & \qw & \ctrl{1} & \qw \\
   \lstick{$\ket0$} & \targ{} & \targ{} & \targ{} & \targ{} & \meter{}\rstick{$Z_1Z_2Z_3Z_4$}
   \end{quantikz}\qquad
   \begin{quantikz}[row sep=0.18cm,column sep=0.3cm]
   \lstick{$d_1$} & \targ{} & \qw & \qw & \qw & \qw \\
   \lstick{$d_2$} & \qw & \targ{} & \qw & \qw & \qw \\
   \lstick{$d_3$} & \qw & \qw & \targ{} & \qw & \qw \\
   \lstick{$d_4$} & \qw & \qw & \qw & \targ{} & \qw \\
   \lstick{$\ket+$} & \ctrl{-4} & \ctrl{-3} & \ctrl{-2} & \ctrl{-1} & \meter{}\rstick{$X_1X_2X_3X_4$}
   \end{quantikz}
\end{center}
The left ancilla ends in $\ket{s}$ with $s$ the parity of the four data bits; the right one, read
in the $X$ basis, gives the parity of their $X$ values. The order of the four CNOTs is not
arbitrary: a CNOT propagates $X$ from control to target and $Z$ from target to control, so an
ancilla error in the middle of the sequence lands on two data qubits (a \emph{hook} error), and
the schedule is chosen so that such pairs lie along the direction in which the code is tolerant.
One cycle is these four layers of CNOTs on every face at once, one measurement and one reset:
on Willow, $\approx1\,\mu$s. Because the ancilla measurement itself errs with probability
$\sim p$, the syndrome is accepted only after it has been seen consistently over $d$ rounds, which
is where the ``$d$ rounds per logical operation'' of \S\ref{ph:sec:ft} comes from.

\subsection{Threshold and the exponential}\label{ph:sec:threshold}

A chain of errors is corrected unless the decoder pairs its ends the wrong way, which it does
only when the chain is at least $d/2$ long. With independent errors at rate $p$, a specific chain
of length $\ell$ has probability $\sim p^\ell$, and the number of chains grows like $A^\ell$ for a
constant $A$ set by the lattice, so the logical error per round is
\[ P_L\ \sim\ \sum_{\ell\ge(d+1)/2}(Ap)^\ell\ \approx\ C\Big(\frac p{p_{\rm th}}\Big)^{(d+1)/2},
   \qquad p_{\rm th}\sim1/A . \]
Below threshold every two steps of $d$ divide $P_L$ by $p_{\rm th}/p$; at threshold the distance
buys nothing and the physical overhead diverges; above it, larger codes are worse. Under
circuit-level depolarising noise with a matching decoder the surface-code threshold is about
$1\%$; the model of \S\ref{rsa:sec:phys} uses $C=0.1$ and $p_{\rm th}=1.22\%$, so that $d=25$ is
the first distance below $10^{-15}$ at $p=10^{-3}$ \cite{ec:gidney25rsa}; the older Fowler fit is
$C=0.03$ with the same exponent \cite{ph:fowler12}. This is also the origin of the engineering
rule that $p$ must be about ten times below threshold: at $p/p_{\rm th}=1/10$ a factor $10$ per
step gives $d\approx27$ for $10^{-15}$; at $1/3$ it would need $d\approx61$ and $7,700$ qubits
per logical qubit; at $1/2$, $d\approx97$. The qubit count is quadratic in $d$ and $d$ is
logarithmic in the target, so the prize is in $p$.

\paragraph{Threshold theorem.} That an arbitrarily long computation can be made reliable at all
is the threshold theorem (Aharonov--Ben-Or, Kitaev, Knill--Laflamme--Zurek, 1996--1999): if the
error per location is below a constant $p_{\rm th}$, with noise that is local and not too
strongly correlated, the logical error can be driven to any target with overhead polylogarithmic
in the circuit size. Early proofs with concatenated codes gave $p_{\rm th}\sim10^{-6}$; the
surface code raised it to $\sim10^{-2}$, which is what made superconducting qubits viable. The
assumptions matter: a cosmic-ray burst or a leaked qubit is neither local nor Pauli, which is why
\S\ref{ph:sec:sc} spent a paragraph on them.

\paragraph{Measured.} Google's Willow experiment (December 2024) ran surface-code memories of
$d=3,5,7$ on 101 qubits and found the logical error per cycle halving at each step
($\Lambda=P_L(d)/P_L(d+2)\approx2.1$), reaching $0.14\%$ per cycle at $d=7$ with a logical
lifetime longer than the best physical qubit's: the first demonstration \emph{below threshold}
\cite{ph:willow24}. A ratio of $2.1$ is far from the $12$ the model assumes at
$p=10^{-3}$; it corresponds to $p/p_{\rm th}\approx1/2$ in the fit above, which is a measure of
how far 2024 hardware was from the hypothesis of the estimates.

\subsection{Decoding, in real time}

The syndrome is measured every cycle because the measurements themselves err, so the decoder sees
a three-dimensional history (two space dimensions, one time) of flipped checks and must pair
them up: minimum-weight perfect matching finds the most probable pairing in time polynomial in
$d$ (empirically $O(d^3)$--$O(d^5)$ per round), union-find does it faster with a small loss in
accuracy, and belief propagation with post-processing serves the qLDPC codes. The result is not
applied as physical gates: it updates a classical \emph{Pauli frame}, and the frame is folded
into the interpretation of later measurements. Two constraints follow. Throughput: the decoder
must consume syndromes at least as fast as the $1\,\mu$s cycle produces them for $\sim10^3$
patches, or a backlog grows without bound. Latency: the frame must be resolved before any
operation whose basis depends on it, which happens at every non-Clifford gate. Google's Willow
decoding ran at $63\,\mu$s latency; FPGA decoders reached sub-microsecond per round on
$d\le12$ codes in 2025--2026; DeepMind's AlphaQubit~2 decodes faster than $1\,\mu$s per cycle
up to $d=11$ and is projected to fail beyond $d\approx23$ \cite{ph:uqt26}. The $10\,\mu$s reaction
time of the model is in this range; the open problem is real-time decoding at $d=25$ for a
thousand patches at once, not decoding as such.

\subsection{Logical gates, and the one that costs}\label{ph:sec:ft}

A code protects a memory. Computing on it requires logical gates that do not spread errors:
a single fault anywhere in the gadget must cause at most one error per code block, so that the
next round of correction removes it. Three constructions do this.

\paragraph{Transversal gates.} Apply the same physical gate to every qubit of the block (or
qubit-by-qubit between two blocks). A fault then touches one qubit per block by construction.
For any CSS code --- a code whose checks are pure $X$ or pure $Z$, like the surface code ---
the transversal CNOT is a logical CNOT: $n$ physical CNOTs between blocks $A$ and $B$ map
$X_A^{\mathbf a}\mapsto X_A^{\mathbf a}X_B^{\mathbf a}$ and $Z_B^{\mathbf b}\mapsto Z_A^{\mathbf b}Z_B^{\mathbf b}$,
so each $X$-check of $A$ becomes itself times the \emph{same} $X$-check of $B$, which is in the
stabiliser; the code space is preserved, and $\bar X_A\mapsto\bar X_A\bar X_B$,
$\bar Z_B\mapsto\bar Z_A\bar Z_B$ is exactly the action of CNOT on logical Paulis. The same
calculation shows why the surface code's transversal $H$ is \emph{not} a logical $H$ (it swaps
$X$- and $Z$-checks, whose layouts differ by a rotation, hence the patch-rotation trick) and why
$T^{\otimes n}$ is never logical for it. The \textbf{Eastin--Knill theorem} (2009) makes this
general: no code that corrects at least one error has a universal set of transversal gates,
because a continuous group of transversal unitaries would commute with the error-detection
projector and let an undetectable error rotate the logical state continuously. Transversal gates
need every qubit of one block to meet its partner in the other, which ions and movable atoms
provide and a fixed chip does not.

\paragraph{Lattice surgery.} On a planar chip the logical CNOT is built from joint
measurements. Two patches side by side have $Z$-type boundaries facing each other. Measuring the
weight-four $Z$-checks that straddle the seam \emph{merges} them into one longer patch; the
product of those new checks is $\bar Z_A\bar Z_B$, so the merge measures the logical operator
$Z_AZ_B$ with outcome $\pm1$ while learning nothing else, since every seam check commutes with
every logical operator of the merged code except that product. Stopping the seam checks
\emph{splits} the patches again. Each merge and split is repeated for $d$ rounds so that
measurement errors on the new checks are corrected, hence the $O(d)$ time per operation and the
$O(d^3)$ space-time volume. A CNOT between control $C$ and target $T$ then uses an ancilla patch
$A$ prepared in $\ket+_L$:
\[ \text{measure }Z_CZ_A\ (m_1),\qquad \text{measure }X_AX_T\ (m_2),\qquad
   \text{measure }Z_A\ (m_3), \]
and apply $X_T^{\,m_1\oplus m_3}$ and $Z_C^{\,m_2}$ as Pauli-frame updates --- a direct calculation on
$\ket\psi_C\ket+_A\ket\phi_T$ (checked numerically for all eight outcome patterns) shows the result
is CNOT$_{C\to T}\ket\psi\ket\phi$. Three joint
measurements of $d$ rounds each, so a logical CNOT costs about $3d$ cycles, $75\,\mu$s at $d=25$,
and needs the ancilla patch plus routing space between $C$ and $T$: that routing is the ``bus''
and the reason the compute region of \S\ref{rsa:sec:phys} reserves patches it never stores
data in. The \textbf{Pauli frame} is the other half of the trick: $X_L$ and $Z_L$ corrections
are never applied physically, only recorded and folded into the interpretation of later
measurements, so Clifford gates ($H$, $S$, CNOT) can even be commuted to the end of the circuit and
absorbed into the measurement bases, leaving the non-Clifford gates as the only operations that
cost time.

\paragraph{What is left.} The surface code thus has cheap Paulis (free), CNOT ($3d$ rounds),
$H$ and $S$ (patch deformation and a cheap injected $\ket{+i}$ state). What it lacks is $T$.

\paragraph{Magic states.} The way round is gate teleportation. Prepare the \emph{magic state}
$\ket T=T\ket+=(\ket0+e^{i\pi/4}\ket1)/\sqrt2$, then
\begin{center}
\begin{quantikz}[row sep=0.35cm,column sep=0.4cm]
   \lstick{$\ket\psi$} & \ctrl{1} & \qw & \gate{S^{m}}\vcw{1} & \qw\rstick{$T\ket\psi$}\\
   \lstick{$\ket T$} & \targ{} & \meter{} & \cw & \cw\rstick{$m$}
   \end{quantikz}
\end{center}
applies $T$ to $\ket\psi$ using only a CNOT, a measurement and an $S$ gate conditioned on the
outcome $m$ --- all Clifford, all cheap --- consuming one $\ket T$. (If $m=1$ the output is $S\,T^\dagger\ket\psi$ up to phase, and $S$ fixes it; the
correction depends on a measurement, so the decoder must have resolved the frame first --- this
is the reaction-time dependence of \S\ref{ph:sec:feeds}.) The difficulty is moved into making
$\ket T$ clean. A noisy $\ket T$, prepared by a physical $T$ gate on one qubit and injected into a
patch, has error $\epsilon\sim p$. \emph{Distillation} (Bravyi--Kitaev 2005) uses a code with a
transversal $T$: the $[[15,1,3]]$ quantum Reed--Muller code has $T^{\otimes15}=\bar T^\dagger$ up
to a Clifford. Encode $\ket+_L$ in it, apply $T$ on all fifteen qubits using fifteen noisy
$\ket T$ states by the teleportation circuit above, measure the code's checks and keep the output
only if no error was detected. A single error among the fifteen is always caught; the output is
wrong only if three or more errors conspire into a logical operator, and counting the
$\binom{15}{3}$ patterns that do gives
\[ \epsilon_{\rm out}\approx35\,\epsilon^3 ,\qquad\text{so}\qquad
   10^{-3}\ \to\ 3.5\times10^{-8}\ \to\ 1.5\times10^{-21}, \]
at the cost of fifteen inputs per output and a success probability $\approx1-15\epsilon$ per
round. Two rounds therefore reach any target the estimates need, and a factory is a few
$[[15,1,3]]$ blocks laid out as surface-code patches, which is where the $O(d^2\log1/\epsilon)$
footprint of magic-state factories comes from. Gidney, Shutty and Jones's \emph{cultivation} (2024)
instead grows a $\ket T$ in a $d=3$ colour code, checks it, grows to $d=5$, checks again, and
converts it to a surface-code patch, post-selecting on every check: about $30,000$ qubit-rounds
per $\ket T$ at $10^{-7}$ error, an order of magnitude cheaper than distillation, which is the
figure \texttt{physical.py} uses \cite{ph:cultivation24}. Cultivation's acceptance rate falls
with the distance of the check, so beyond about $10^9$ $T$ states per run it is used to seed
distillation rather than replace it \cite{ph:uqt26}.

\paragraph{The 25 microseconds.} A Toffoli is a CCZ between Hadamards, and a CCZ is made from
four $\ket T$ states with the temporary-AND construction of Part~V (\S\ref{sec:tempand}), or
eight in the direct form the factory uses. The model's factory is a block of $3\times4$ patches
of distance $25$, $12\times1,352=16,224$ qubits. Cultivating eight $\ket T$ states costs
$8\times30,000=240,000$ qubit-rounds, that is $240,000/16,224\approx15$ rounds of the whole
factory; assembling them into a CCZ state takes six layers of lattice surgery of $2d/3$ rounds
each, $100$ rounds; the total $115$ is padded to $150$ for slack. Six factories in parallel
therefore deliver one CCZ every $150/6=25$ rounds, and a round is the $1\,\mu$s cycle: one
Toffoli every $25\,\mu$s, which is \S\ref{rsa:sec:phys}'s ``counting time'' and the origin of
every runtime in Part~VIII. The compute region of $7\times18$ patches, $170,352$ qubits, holds
the six factories and the lattice-surgery workspace that consumes their output. Since $25$ rounds
is also one logical cycle at $d=25$, the machine is factory-bound and logic-bound in equal
measure; a seventh factory would not help without more workspace.

\subsection{Beyond the surface code}

The surface code pays $2d^2$ qubits per logical qubit because every check is local on a plane:
the Bravyi--Poulin--Terhal bound $kd^2=O(n)$ makes that the best any two-dimensional local code
can do. Quantum LDPC codes keep the checks sparse but let them reach across the chip, and then
$k$ and $d$ can both grow linearly with $n$ (Panteleev--Kalachev, 2021). IBM's \emph{gross} code
$[[144,12,12]]$ stores twelve logical qubits in $144$ data qubits with weight-6 checks, needing
degree-6 connectivity, which is what the Loon couplers provide; Iceberg Quantum's Pinnacle
architecture (2026) stacks such codes into an RSA-2048 machine of under $10^5$ physical qubits
at the price of a month's runtime and degree-8 to 10 connectivity, which Google argued is not
feasible on superconducting chips \cite{ph:pinnacle26,ph:uqt26}. Google's answer stays planar:
\emph{yoked} surface codes (2023) halve the qubits per logical qubit by adding a few parity
checks across patches, used for RSA's idle exponent register in \S\ref{rsa:sec:phys} at
$430$ qubits per logical qubit, and a ``denser planar surface code'' (May 2026) claims a further
factor $4.5$ \cite{ph:uqt26}. On atoms and ions, where qubits move, the connectivity is free
and the qLDPC codes are the natural choice, which is why the 2026 atom-based estimates
(\S\ref{ph:sec:landscape}) use far fewer qubits and far more time.

\begin{center}
\begin{tikzpicture}[x=112mm,y=38mm]
  \fill[okfill] (0.0005,1) plot[domain=0.0005:0.1893,samples=60] (\x,{1-\x*ln(3)/ln(2)+\x*ln(\x)/ln(2)+(1-\x)*ln(1-\x)/ln(2)}) -- (0,0) -- cycle;
  \fill[lkfill] (0.0005,1) plot[domain=0.0005:0.1893,samples=60] (\x,{1-\x*ln(3)/ln(2)+\x*ln(\x)/ln(2)+(1-\x)*ln(1-\x)/ln(2)}) -- (0.3786,0) plot[domain=0.3786:0.0005,samples=60] (\x,{1-(\x/2)*ln(3)/ln(2)+(\x/2)*ln(\x/2)/ln(2)+(1-\x/2)*ln(1-\x/2)/ln(2)}) -- cycle;
  \draw[->] (0,0) -- (0.53,0); \node[font=\scriptsize] at (0.27,-0.14) {relative distance $\delta=d/n$};
  \draw[->] (0,0) -- (0,1.08) node[above right,font=\scriptsize] {rate $k/n$};
  \foreach \t in {0.1,0.2,0.3,0.4,0.5} \draw (\t,0) -- (\t,-0.02) node[below,font=\tiny] {\t};
  \foreach \t in {0.5,1} \draw (0,\t) -- (-0.01,\t) node[left,font=\tiny] {\t};
  \draw[cverm,thick] (0,1) -- (0.5,0);
  \draw[cpurple,thick,domain=0.0005:0.3786,samples=70] plot (\x,{1-(\x/2)*ln(3)/ln(2)+(\x/2)*ln(\x/2)/ln(2)+(1-\x/2)*ln(1-\x/2)/ln(2)});
  \draw[cblue,thick,domain=0.0005:0.1893,samples=70] plot (\x,{1-\x*ln(3)/ln(2)+\x*ln(\x)/ln(2)+(1-\x)*ln(1-\x)/ln(2)});
  \draw[cblue,thick,dashed,domain=0.0005:0.1100,samples=50] plot (\x,{1+2*\x*ln(\x)/ln(2)+2*(1-\x)*ln(1-\x)/ln(2)});
  \node[font=\scriptsize,color=cverm,anchor=west] at (0.37,0.33) {Singleton};
  \node[font=\scriptsize,color=cpurple,anchor=west] at (0.215,0.42) {Hamming};
  \node[font=\scriptsize,color=cblue,anchor=west] at (0.13,0.29) {GV};
  \node[font=\scriptsize,color=cblue,anchor=west] at (0.108,0.06) {CSS};
  \fill[black] (0.004,0.004) circle (1.2pt) node[above right,font=\tiny,xshift=1pt] {surface};
\end{tikzpicture}
\end{center}
{\small\noindent Which large codes can exist, in rate $k/n$ against relative distance $\delta=d/n$.
Gilbert--Varshamov: codes exist with $k/n\ge1-\delta\log_23-H_2(\delta)$ (green), CSS codes with
$1-2H_2(\delta)$ (dashed). Singleton: $k/n\le1-2\delta$ for every code. Hamming:
$k/n\le1-\frac\delta2\log_23-H_2(\frac\delta2)$ is proven for non-degenerate codes and for
$\delta\le0.235$ or $d=3$ in general; whether degenerate codes can beat it elsewhere is open, so
the peach band is ``unknown'' above GV and ``impossible'' only above Singleton. The surface code
sits at the origin: vanishing rate and vanishing relative distance, the price of local checks.}

\subsection{Logical qubits, as demonstrated}

\begin{center}\footnotesize
\renewcommand{\arraystretch}{1.12}
\begin{tabular}{@{}llp{2.8cm}lp{5.2cm}@{}}
\toprule
when & who & code & physical / logical & what was shown\\
\midrule
2021 & Maryland & Bacon--Shor-13, ions & 15 / 1 & first fault-tolerant logical qubit; logical fidelity $99.7\%$\\
2022 & Quantinuum & Steane colour, ions & 10 / 1 & logical SPAM better than physical (break-even)\\
2022 & Google & surface $d=5$, SC & 49 / 1 & logical error \emph{not} better than $d=3$: above threshold\\
Dec 2023 & Harvard--QuEra & colour + qLDPC, atoms & 280 / 48 & logical processor, transversal gates, CCZ; post-selected\\
Dec 2024 & Google & surface $d=7$, SC & 101 / 1 & \textbf{below threshold}: $\Lambda\approx2.1$, $0.14\%$ per cycle\\
2024 & Quantinuum--Microsoft & tesseract / Steane, ions & 56 / 12 & 12-qubit logical GHZ states, post-selected\\
2025 & QuEra & colour + qLDPC, atoms & 200 of 448 & logical CZ by surgery and teleportation, CCZ, rotations; atom-loss post-selection\\
2026 & Quantinuum & Steane, Helios ions & 98 / 48 & $H$, CZ, $T$ on 48 logical qubits; teleported $T$ states; small QAOA and HHL\\
2026 & China (Hefei et al.) & surface $d=7$, SC & 125 / 2 & two $d=7$ patches, transversal teleportation\\
\bottomrule
\end{tabular}
\end{center}
\noindent From \cite{ph:uqt26} (Tables 9--10). None is a full fault-tolerant stack: as of
September 2026 no experiment had run a distilled or cultivated $T$ state through a below-threshold
memory with real-time decoding and no post-selection, and the largest code distance on hardware
was $7$ against the $25$ the estimates need.

% ===========================================================================
\section{Where things stand}\label{ph:sec:landscape}

\subsection{Milestones}

\begin{center}\small
\renewcommand{\arraystretch}{1.1}
\begin{tabular}{@{}lp{12.4cm}@{}}
\toprule
1994 & Shor's factoring and discrete-logarithm algorithms. \\
1995 & Cirac--Zoller trapped-ion gate proposal; first BEC. \\
1996--97 & Shor and Steane codes; Gottesman's stabiliser formalism; Kitaev's toric code; threshold theorem. \\
1999 & Nakamura: first coherent superconducting (charge) qubit, $\sim1$\,ns coherence. \\
2001 & Shor's algorithm factors 15 on a 7-qubit NMR machine (compiled). KLM linear optics. \\
2007 & The transmon. \\
2010 & First Rydberg-blockade CNOT between neutral atoms. \\
2012 & Fowler et al.: surface-code architecture with $\sim1\%$ threshold; 21 factored photonically (compiled). \\
2016 & Monz et al.: scalable Shor circuit factors 15 with trapped ions. \\
2015 & Loophole-free Bell tests (NV centres at Delft; photons at Vienna and NIST). \\
2019 & Sycamore random-circuit sampling (``supremacy''); Gidney--Eker\aa: RSA-2048 in 8 hours with 20 million qubits. \\
2021 & Maryland: first fault-tolerant logical qubit (ions); Gouzien--Sangouard memory-based estimate; good qLDPC codes (Panteleev--Kalachev). \\
2023 & Harvard--QuEra 48 logical qubits; Evered: $99.5\%$ Rydberg CZ on 60 atoms; Litinski: ECDLP-256 in 50 million Toffolis. \\
2024 & Google Willow below threshold; Caltech $6,100$-atom array; Gidney--Shutty--Jones cultivation; IBM gross code. \\
2025 & Gidney: RSA-2048 under a million qubits in under a week; Quantinuum Helios 98 ions at $99.92\%$; Oxford $1.5\times10^{-7}$ single-qubit error; QuEra fault-tolerant architecture paper. \\
2026 & Pinnacle: RSA-2048 under $10^5$ qubits (qLDPC); IonQ: ECDLP-256 in 26 days on $20,000$ ions; Quantinuum 48 logical qubits with $T$; Nighthawk and Loon; the ECDLP circuits of Part~VII. \\
\bottomrule
\end{tabular}
\end{center}

\subsection{What has actually been factored}

Every Shor demonstration to date has been compiled for its answer. The 2001 NMR and 2016 ion
experiments factored $15$; the 2012 photonic experiment factored $21$ with qubit recycling; and
$21$ is still the largest integer factored by a Shor-type circuit on hardware in 2026. The reason
is in Part~V: a circuit for $15$ collapses to $21$ entangling gates because $a^x\bmod15$ has a
trivial structure, while a general circuit for $21$ needs $2,405$ entangling gates
(Gidney, 2025), beyond what any uncorrected processor runs coherently
\cite{ph:uqt26}. Larger ``factorings'' --- $35$ on IBM hardware, $56,153$, $1,099,551,473,989$
with variational methods, $2048$-bit integers with specially close factors on an annealer, the
2026 ``Q-Day prize'' for a 15-bit elliptic-curve key whose $98,000$-gate circuit on 65 qubits
succeeded with probability $10^{-215}$ and was filtered by a classical checker --- are not Shor's
algorithm and do not scale; the $2022$ claim of RSA-2048 with $372$ qubits (Schnorr's lattice
method plus QAOA) was shown not to scale past $70$ bits even with perfect qubits. Classically,
the record is RSA-250, $829$ bits, with the number field sieve in 2020 at about $2,700$
core-years; a $862$-bit RSA-260 was claimed in September 2026 without an audited computation
\cite{ph:uqt26}.

\subsection{The resource-estimate landscape}

\begin{center}\footnotesize
\renewcommand{\arraystretch}{1.12}
\begin{tabular}{@{}p{3.1cm}p{2.1cm}p{2.5cm}p{1.7cm}p{5.0cm}@{}}
\toprule
estimate & target & physical qubits & time & assumptions and caveats\\
\midrule
Gidney--Eker\aa\ 2019 \cite{gidney-ekera} & RSA-2048 & $20$\,M & 8 h & surface code, $p=10^{-3}$, $1\,\mu$s cycle\\
Gouzien--Sangouard 2021 & RSA-2048 & $13,436$ + memory & 177 d & $2$ logical qubits plus a $28$-million-mode quantum memory\\
Alice \& Bob 2019 / 2023 \cite{ph:cats19,ph:cats23} & RSA-2048 / ECC-256 & $349$\,k / $126$\,k cat qubits & 4 d / 9 h & repetition code on biased-noise cats\\
Gidney 2025 \cite{ec:gidney25rsa} & RSA-2048 & $<1$\,M & $<1$ week & residue arithmetic, yoked codes, cultivation; the model of Part VIII\\
Pinnacle 2026 \cite{ph:pinnacle26} & RSA-2048 & $94$\,k / $135$\,k / $381$\,k & 1 mo / 1 wk / 1 d & qLDPC, degree-8 to 10 connectivity, SC\\
IonQ 2026 \cite{ec:ionq26} & ECC-256 & $20,000$ & 26 d & walking-cat codes, $99.99\%$ gates, $63\%$ success\\
Oratomic 2026 & RSA-2048 / ECC-256 & $10^4$--$10^5$ & months--years & qLDPC on atoms; slow cycle\\
Chicago 2026 (distributed) & RSA-2048 & $\sim$M & $\sim1$ yr & modular SC with interconnects\\
this tutorial, Part VIII & ECC-256 / RSA-2048 & $2.14$\,M / $898$\,k & 15 min / 5 d & Gidney 2025's model; circuits of Parts VII--VIII\\
\bottomrule
\end{tabular}
\end{center}
\noindent Compiled from \cite{ph:uqt26} (Figures 285--286) and the papers cited; Google's 2026
elliptic-curve estimate \cite{ec:babbush26}, the surface-code reference for ECC-256, is compared
with Part~VII's circuits there rather than here. All assume
$99.9\%$ fidelities ($99.99\%$ for IonQ) uniformly at scale, independent Pauli noise, and a
monolithic processor unless stated. The trend from $20$ million qubits to under $10^5$ in seven
years came from arithmetic (Parts V--VIII), from codes (yoked, qLDPC), from magic states
(cultivation) and from trading time for space; it did not come from hardware, which in the same
period went from no logical qubit to one below-threshold memory at $d=7$.

\subsection{Forecasts}

Mosca's inequality frames the question for a defender: data must stay secret for $D$ years,
migration to post-quantum cryptography takes $T$ years, and a cryptographically relevant
quantum computer arrives in $G_c$ years; if $D+T>G_c$ the migration is already late. $D$ and $T$
are knowable; $G_c$ is not. A 2020 statistical study of expert forecasts gave a $90\%$ interval of
2039--2058 with a median of 2050 for a Shor attack on RSA-2048, and 2026--2033 for a
proof-of-concept fault-tolerant machine; the yearly evolutionQ expert surveys give similar
spreads; Shor's own guess was ``between 20 and 40 years'' with the possibility of never. Vendor
roadmaps are more aggressive --- IBM's Starling (200 logical qubits, $10^8$ gates) for 2029 and
Blue Jay ($2,000$ logical qubits, $10^9$ gates) for 2033, Quantinuum's universal fault-tolerant
Apollo for 2029, Microsoft's fault-tolerant machine for 2029 --- and past NISQ roadmaps slipped
by one to three years \cite{ph:uqt26}. Against the numbers of this tutorial: ECDLP-256 needs
about $1,600$ logical qubits running $3.7\times10^7$ Toffolis with a $10^{-15}$ error per cycle,
which is a Blue Jay-class machine with a hundredfold lower logical error than Starling's target
and a $25\,\mu$s Toffoli; none of the 2026 demonstrations is within four orders of magnitude of it
in any single dimension.

\begin{keybox}[title={The physics, in one line each}]
A qubit is a two-level system with a Hamiltonian we can switch; a gate is a rotation, produced by
a resonant drive whose phase picks the axis and whose area picks the angle; an entangling gate is
an interaction switched on for a time, and it is the slowest and dirtiest operation on every
platform. An error is a channel on the density operator, and a stabiliser measurement projects
its continuum onto Paulis with a syndrome, which is why a code can correct it. Below a threshold
near $1\%$ the surface code's logical error falls as $(p/p_{\rm th})^{(d+1)/2}$, so the physical
error rate is the one number that matters exponentially. A $T$ gate is teleported from a magic
state that a factory makes in $\sim150$ rounds, and six factories give one Toffoli per
$25\,\mu$s on a $1\,\mu$s-cycle machine. In 2026 the hardware had one below-threshold memory at
$d=7$ and $98$ ions at $99.92\%$; the estimates need $d=25$ on a thousand patches.
\end{keybox}
